\color{unidarkgreen}{\subsection{Аварийная максимальная токовая защита}}
\color{black}

\begin{enumerate}[label=\arabic{section}.\arabic{subsection}.\arabic*,  align=left, labelsep=4pt, leftmargin=0pt, itemindent=57pt]

\item
Аварийная максимальная токовая защита (АМТЗ) применяется в качестве аварийной защиты при
неисправностях в цепях напряжения.

\item
В составе АМТЗ предусмотрены:
\begin{itemize}
\item 2 ступени, реагирующие на величину фазных токов (АМТЗ-Ф);
\item 2 ступени, реагирующие на величину тока обратной последовательности (АМТЗ-ОП);
\item 2 ступени, реагирующие на величину тока нулевой последовательности (АМТЗ-НП).
\end{itemize}

\begin{figure}[H]
\centering
\includegraphics[width=1\textwidth,height=1\textheight,keepaspectratio]{\fbpath/661.0161 АМТЗ/img/AMTZ.pdf}
\captionof{figure}{Функциональная схема ступеней АМТЗ-Ф}\label{pict_amtz:AMTZ}
\end{figure}

\item
Пороги тока срабатывания ступеней АМТЗ, реагирующих на фазный ток, задаются уставками \textbf{«Iср АМТЗ-
Ф-1»} и \textbf{«Iср АМТЗ-Ф-2»}. Действие ступеней осуществляется с регулируемыми независимыми выдержками
времени \textbf{«Тср АМТЗ-Ф-1»} и \textbf{«Тср АМТЗ-Ф-2»}. Ввод в работу осуществляется с помощью программных ключей
\textbf{«АМТЗ-Ф-1»} и \textbf{«АМТЗ-Ф-2»}.

\item
При необходимости можно ввести контроль от БК-I-м отдельно для каждой ступени с помощью
программных переключателей \textbf{«БК-I-м АМТЗ-Ф-1»} и \textbf{«БК-I-м АМТЗ-Ф-2»}.

\item
Ввод в работу ступеней, реагирующих на ток обратной последовательности АМТЗ-ОП и на ток нулевой
последовательности АМТЗ-НП, настройка параметров срабатывания, ввод контроля от БК-I-м осуществляется
аналогично.

\item
Автоматический ввод в работу АМТЗ при возникновении неисправности в цепях напряжения от ТН 1СШ
производится при отсутствии внешнего сигнала \textit{«Вывод АМТЗ»}. Пуски ступеней возможны даже при отсутствии
неисправности в цепях напряжения, но выходная логика срабатывания находится в заблокированном
состоянии. При появлении входного сигнала \textit{«БНН»} выходная логика ступеней на срабатывание
деблокируется, при этом формируется логический сигнал \textit{«АМТЗ активна»}. Предусмотрена также
возможность оперативного ввода АМТЗ (независимо от сигнала срабатывания БНН) от команды управления 
\textit{«Опер. ввод АМТЗ»}.

\item
Оперативный ввод/вывод АМТЗ производится с помощью БУ (блоков управления функциями) \textit{«Вывод АМТЗ»}. 

\begin{figure}[H]
\centering
\includegraphics[width=1\textwidth,height=1\textheight,keepaspectratio]{\fbpath/661.0161 АМТЗ/img/GL_AMTZ.pdf}
\captionof{figure}{Функциональная схема общей логики АМТЗ}\label{pict_amtz:GL_AMTZ}
\end{figure}

\small
\setlength{\extrarowheight}{0.05cm} % высота ячеек в таблице
\begin{longtable}
{|>{\centering\arraybackslash}m{4.5cm}|>{\centering\arraybackslash}m{4.5cm}|>{\centering\arraybackslash}m{3cm}|>{\centering\arraybackslash}m{3.5cm}|}
\caption{Перечень уставочных параметров АМТЗ\hfill\vspace{-0.5\baselineskip}}\label{amtz:tbl1} \\ 
\hhline{|-|-|-|-|}
\rowcolor{gray!30}
Наименование уставки & Диапазон & Шаг &  Значение по умолчанию \\ 
\hhline{|-|-|-|-|}
\endfirsthead

\caption*{\hspace{3pt}\emph{Продолжение таблицы \ref{amtz:tbl1}\hfill\vspace{-0.5\baselineskip}}} \\ % сделано по ГОСТ 2.105 п.6.8.7
\hhline{|-|-|-|-|}
\rowcolor{gray!30}
Наименование уставки & Диапазон & Шаг &  Значение по умолчанию \\  
\hhline{|-|-|-|-|}
\endhead
\hhline{|-|-|-|-|}

\endfoot
\endlastfoot

%===m>АМТЗ|-
%===m

\end{longtable}
\setlength{\extrarowheight}{0.0cm} % возврат параметра высота ячеек в таблице
\normalsize





























































\end{enumerate} 